\section{Literature Survey}\label{sec:literature}

\textit{Rebalancing policy.} Portfolio construction selects weights that trade
expected return against variance~\cite{markowitz1952}, and risk-adjusted
performance measures make the resulting drift consequential~\cite{sharpe1966}.
Practice answers drift with calendar or threshold rules, neither free of
opportunity cost since restoring target weights trims positions whose
intermediate-horizon returns are reported to continue rather than
revert~\cite{jegadeesh1993}. The applied literature that proposes machine
learning here argues the point conceptually: Adedoyin catalogues the model
families that might replace fixed-interval rebalancing and the overfitting
risk they face~\cite{adedoyin_ai}, and Kumari reaches a similar conclusion for
threshold policies, naming logistic regression as the appropriate device for a
binary financial outcome~\cite{kumari2024}. Both are reviews; neither states a
reproducible method, a labelling procedure or a numeric result, and this work
uses them only as motivation and as precedent for the model class.

\textit{Financial ML and federated mechanics.} Mashrur \emph{et al.} survey
machine learning for financial risk management and establish supervised
classifiers, logistic regression among them, as standard instruments for
binary financial determinations, while also naming generated data and
federated learning as directions for settings where confidential or rare-event
data are unavailable~\cite{mashrur2020}. Deep architectures apply the same
representation-learning spirit to asset universes~\cite{heaton2017}. Neither
models the portfolio-level rebalancing action itself. Federated averaging
trains a shared model by repeatedly averaging local updates instead of
centralising data~\cite{mcmahan2017}, applicable across the horizontal,
vertical and transfer settings surveyed by Yang \emph{et al.}~\cite{yang2019federated};
its dominant obstacle is statistical heterogeneity, which degrades convergence
under non-independent client data~\cite{li2020challenges,li2020noniid}, and
surveys of federated recommendation add that retaining data on-device removes
a class of collection risk while introducing inference and poisoning threats
of its own~\cite{javeed2024}.

\textit{Financial federated evidence.} Noseda \emph{et al.} train a shared
volatility classifier and report federated averaging reaching accuracy close
to centralised training under an IID partition ($0.7162$ vs.\ $0.7295$) and
exceeding isolated local training under a non-IID partition ($0.7206$ vs.\
$0.5366$), while strong differential-privacy noise collapses utility below
every non-private configuration~\cite{noseda2025}. Yu, Wynter and Lim apply
federated reinforcement learning to portfolio management and find that plain
aggregation causes training to fail, motivating a personalised, regularised
server objective~\cite{yu2022fedrl}. The lesson is that collaboration can
approach centralised quality on heterogeneous financial data, but neither
privacy mechanisms nor plain averaging are free of cost.

Table~\ref{tab:t14} positions the present work against this survey. No
identified work treats whether a drifted portfolio requires rebalancing now as
supervised binary classification trained on forward, net-of-cost outcome
labels, and none measures the degeneracy that arises when a rebalancing model
is trained on the labels of the rule it replaces. This paper occupies that
slot: it formulates rebalancing-necessity detection on backtested outcome
labels, quantifies the resulting degeneracy against heuristic labels, and
couples the detector to a federated design exchanging model updates rather
than portfolio records.
